\section{Related work}
\label{sec:related}

\subsection{Genetic and multi-objective evolutionary fuzzy systems}
Genetic fuzzy systems (GFSs) use evolutionary search to tune or learn the data base
(membership functions) and/or the rule base of a fuzzy system~\cite{herrera2008}.
Ishibuchi et al.\ formulated rule selection as a single- or two-objective problem,
trading classification accuracy against the number of rules~\cite{ishibuchi1997}, and
later analysed the interpretability--accuracy trade-off with multi-objective fuzzy
genetics-based machine learning~\cite{ishibuchi2007}. Alcal\'a et al.\ learned rule and
data bases concurrently with a multi-objective evolutionary algorithm~\cite{alcala2009},
and Antonelli et al.\ added the \emph{integrity} of the fuzzy partitions as a third
objective alongside accuracy and complexity~\cite{antonelli2011}. Surveys of
interpretability measures~\cite{gacto2011} and of MOEFS
applications~\cite{fazzolari2013} show that complexity-based measures (number of
rules, number of conditions) dominate in practice, and that the accuracy objective is
almost always a classification-error criterion. Evolutionary fuzzy systems have since
been positioned as a route to explainable AI~\cite{fernandez2019}. The NSGA-II
algorithm~\cite{deb2002}, with fast non-dominated sorting and crowding distance, and the
hypervolume indicator~\cite{zitzler1999} are standard tools in this line of work.

\subsection{Fuzzy and genetic-fuzzy models in medicine}
Pe\~na-Reyes and Sipper evolved compact fuzzy rule bases for breast-cancer
diagnosis~\cite{penareyes1999}. Nauck and Kruse obtained interpretable fuzzy
classification rules from medical data with neuro-fuzzy learning~\cite{nauck1999}.
Hybrid neuro-fuzzy and fuzzy systems have been applied to diabetes and heart
disease~\cite{kahramanli2008,anooj2012}, and hybrid intelligent systems combining fuzzy
min--max networks with decision trees and random forests to several medical
benchmarks~\cite{seera2014}. Gorza{\l}czany and Rudzi\'nski proposed a multi-objective
genetic-fuzzy approach for interpretable and accurate medical data
classification~\cite{gorzalczany2017}, and Pota et al.\ discussed design principles for
rule-based fuzzy classifiers in medicine~\cite{pota2017}. ANFIS~\cite{jang1993}
popularised \emph{hybrid learning}, in which antecedent and consequent parameters are
estimated by different procedures. We use the same idea, combining evolutionary
search over antecedents with convex estimation of consequents. In all of these works
the probabilistic output, when there is one, is not optimised or assessed for calibration.

\subsection{Probability calibration}
A forecaster is calibrated if, among cases assigned probability $p$, the event occurs
with frequency $p$. The Brier score~\cite{brier1950} and the logarithmic score are
strictly proper scoring rules: their expectation is uniquely minimised by the true
probabilities, and they reward calibration and sharpness
together~\cite{gneiting2007}. Post-hoc calibration maps scores to probabilities with
Platt scaling or isotonic regression~\cite{zadrozny2002,niculescu2005}. Binned
calibration errors (ECE) are widely reported~\cite{naeini2015,guo2017} but depend on
the binning; equal-mass binning is one remedy~\cite{nixon2019}. In clinical
prediction, the calibration intercept and slope from a logistic recalibration
model are recommended, together with the calibration
hierarchy~\cite{steyerberg2010,vancalster2016,vancalster2019,huang2020}.

\subsection{Benchmarking on tabular clinical data}
Tree-based ensembles still outperform deep tabular models on medium-sized
data~\cite{grinsztajn2022,shwartzziv2022}. The main implementations are
XGBoost~\cite{chen2016}, LightGBM~\cite{ke2017} and CatBoost~\cite{prokhorenkova2018}. In
clinical prediction a systematic review found no performance benefit of machine
learning over logistic regression~\cite{christodoulou2019}, so a well-tuned logistic
model is an essential baseline. Fair comparisons require nested model selection to
avoid optimistic bias~\cite{cawley2010}, strict separation of fitted preprocessing from
test data~\cite{kaufman2012}, and appropriate statistics across
datasets~\cite{demsar2006} and across resampled folds~\cite{nadeau2003}. PR curves add
information on imbalanced data~\cite{saito2015}, and reporting should follow
TRIPOD~\cite{collins2015}.
